% ******************************************************************************
% ****************************** Custom Margin *********************************

% Add `custommargin' in the document class options to use this section
% Set {innerside margin / outerside margin / topmargin / bottom margin}  and
% other page dimensions
\ifsetCustomMargin
  \RequirePackage[left=37mm,right=30mm,top=35mm,bottom=30mm]{geometry}
  \setFancyHdr % To apply fancy header after geometry package is loaded
\fi

% Add spaces between paragraphs
%\setlength{\parskip}{0.5em}
% Ragged bottom avoids extra whitespaces between paragraphs
\raggedbottom
% To remove the excess top spacing for enumeration, list and description
%\usepackage{enumitem}
%\setlist[enumerate,itemize,description]{topsep=0em}

% *****************************************************************************
% ******************* Fonts (like different typewriter fonts etc.)*************

% Add `customfont' in the document class option to use this section

\ifsetCustomFont
  % Set your custom font here and use `customfont' in options. Leave empty to
  % load computer modern font (default LaTeX font).
  %\RequirePackage{helvet}

  % For use with XeLaTeX
  %  \setmainfont[
  %    Path              = ./libertine/opentype/,
  %    Extension         = .otf,
  %    UprightFont = LinLibertine_R,
  %    BoldFont = LinLibertine_RZ, % Linux Libertine O Regular Semibold
  %    ItalicFont = LinLibertine_RI,
  %    BoldItalicFont = LinLibertine_RZI, % Linux Libertine O Regular Semibold Italic
  %  ]
  %  {libertine}
  %  % load font from system font
  %  \newfontfamily\libertinesystemfont{Linux Libertine O}
\fi

% *****************************************************************************
% **************************** Custom Packages ********************************

% ************************* Algorithms and Pseudocode **************************

% The algorithmicx / algorithms bundle is not present in all TeX installations
% (e.g. texlive-basic), so the thesis defines a lightweight, self-contained
% pseudocode facility from `float` and `enumitem` only, both of which are always
% available. It provides: an `algorithm` float (numbered, caption above, listed
% via \listofalgorithms if desired) and a small set of pseudocode macros used in
% Appendix~\ref{app:pseudocode}.
\usepackage{float}
\usepackage{enumitem}

% A floating environment for algorithms, styled with rules (caption on top).
\floatstyle{ruled}
\newfloat{algorithm}{htbp}{loa}[chapter]
\floatname{algorithm}{Algorithm}

% Pseudocode body: a tight, auto-numbered line list emulating the algorithmicx
% `algorithmic' environment. The optional [n] argument (line-number stride) is
% accepted and ignored; every line is numbered.
\newlist{algcode}{enumerate}{1}
\setlist[algcode]{label=\footnotesize\arabic*:, ref=\arabic*,
  leftmargin=2.8em, itemsep=1pt, topsep=3pt, parsep=0pt}
\newenvironment{algorithmic}[1][0]{\footnotesize\begin{algcode}}{\end{algcode}}
% Structural keywords and helpers matching the algorithmicx names used in the text.
\newcommand{\algkw}[1]{\textbf{#1}}
\newcommand{\Require}[1]{\item[\textbf{Input:}] #1}
\newcommand{\Ensure}[1]{\item[\textbf{Output:}] #1}
\newcommand{\State}{\item}
\newcommand{\Comment}[1]{{\normalfont\itshape\quad// #1}}
\newcommand{\For}[1]{\item \algkw{for} #1 \algkw{do}}
\newcommand{\EndFor}{\item \algkw{end for}}
\newcommand{\While}[1]{\item \algkw{while} #1 \algkw{do}}
\newcommand{\EndWhile}{\item \algkw{end while}}
\newcommand{\If}[1]{\item \algkw{if} #1 \algkw{then}}
\newcommand{\EndIf}{\item \algkw{end if}}


% ********************Captions and Hyperreferencing / URL **********************

% Captions: This makes captions of figures use a boldfaced small font.
%\RequirePackage[small,bf]{caption}

\RequirePackage[labelsep=space,tableposition=top]{caption}
\renewcommand{\figurename}{Fig.} %to support older versions of captions.sty


% *************************** Graphics and figures *****************************

%\usepackage{rotating}
%\usepackage{wrapfig}

% Uncomment the following two lines to force Latex to place the figure.
% Use [H] when including graphics. Note 'H' instead of 'h'
%\usepackage{float}
%\restylefloat{figure}

% \FloatBarrier prevents floats from drifting past section boundaries.
\usepackage{placeins}

% ****************************** Vector schematics ******************************
% TikZ supplies the apparatus and process schematics of Chapter~\ref{ch:fab}, which
% are drawn inline rather than imported as bitmaps so that their type matches the
% body text. Only libraries present in texlive-basic are loaded: `arrows.meta' is
% NOT available in this installation, so the classic `arrows' tips are used.
\usepackage{tikz}
\usetikzlibrary{arrows,positioning,calc,fit,patterns,shapes.geometric,shapes.misc,
  decorations.pathmorphing,backgrounds}

% Float placement: relax the LaTeX defaults so large figures share a page with
% text instead of being exiled to half-empty float-only pages. Raise the share
% of a page that floats may occupy, require float pages to be nearly full, and
% allow more floats per page.
\renewcommand{\topfraction}{0.9}
\renewcommand{\bottomfraction}{0.8}
\renewcommand{\textfraction}{0.07}
\renewcommand{\floatpagefraction}{0.85}
\setcounter{topnumber}{3}
\setcounter{bottomnumber}{2}
\setcounter{totalnumber}{4}

% Subcaption package is also available in the sty folder you can use that by
% uncommenting the following line
% This is for people stuck with older versions of texlive
%\usepackage{sty/caption/subcaption}
\usepackage{subcaption}

% ********************************** Tables ************************************
\usepackage{booktabs} % For professional looking tables
\usepackage{multirow}

%\usepackage{multicol}
\usepackage{longtable}
\usepackage{tabularx}
% Centred equal-width column for the source-data tables of Appendix~\ref{app:sourcedata}.
\newcolumntype{Y}{>{\centering\arraybackslash}X}


% *********************************** SI Units *********************************
\usepackage{siunitx} % use this package module for SI units

% ******************************* Code Listings *******************************
% Verbatim source listings for Appendix C (the reproducible code snippet).
\usepackage{xcolor}
\usepackage{listings}
\definecolor{lstbg}{gray}{0.97}
\definecolor{lstcomment}{rgb}{0.30,0.45,0.30}
\definecolor{lstkw}{rgb}{0.20,0.20,0.60}
\definecolor{lststr}{rgb}{0.55,0.25,0.10}
\lstdefinestyle{thesiscode}{%
  language=Python,
  basicstyle=\ttfamily\scriptsize,
  backgroundcolor=\color{lstbg},
  keywordstyle=\color{lstkw},
  commentstyle=\color{lstcomment}\itshape,
  stringstyle=\color{lststr},
  numbers=left, numberstyle=\tiny\color{gray}, numbersep=6pt,
  showstringspaces=false,
  breaklines=true, breakatwhitespace=false,
  postbreak=\mbox{\textcolor{gray}{$\hookrightarrow$}\space},
  tabsize=4, columns=fullflexible, keepspaces=true,
  frame=single, rulecolor=\color{gray!40},
  captionpos=b, xleftmargin=1.5em, framexleftmargin=1.5em,
}
\lstset{style=thesiscode}


% ******************************* Line Spacing *********************************

% Choose linespacing as appropriate. Default is one-half line spacing as per the
% University guidelines

% \doublespacing
% \onehalfspacing
% \singlespacing


% ************************ Formatting / Footnote *******************************

% Don't break enumeration (etc.) across pages in an ugly manner (default 10000)
%\clubpenalty=500
%\widowpenalty=500

%\usepackage[perpage]{footmisc} %Range of footnote options


% *****************************************************************************
% *************************** Bibliography  and References ********************

%\usepackage{cleveref} %Referencing without need to explicitly state fig /table

% Add `custombib' in the document class option to use this section
\ifuseCustomBib
   \RequirePackage[square, sort, numbers, authoryear]{natbib} % CustomBib

% If you would like to use biblatex for your reference management, as opposed to the default `natbibpackage` pass the option `custombib` in the document class. Comment out the previous line to make sure you don't load the natbib package. Uncomment the following lines and specify the location of references.bib file

%\RequirePackage[backend=biber, style=numeric-comp, citestyle=numeric, sorting=nty, natbib=true]{biblatex}
%\bibliography{References/references} %Location of references.bib only for biblatex

\fi

% changes the default name `Bibliography` -> `References'
\renewcommand{\bibname}{References}


% ******************************************************************************
% ************************* User Defined Commands ******************************
% ******************************************************************************

% *********** To change the name of Table of Contents / LOF and LOT ************

%\renewcommand{\contentsname}{My Table of Contents}
%\renewcommand{\listfigurename}{My List of Figures}
%\renewcommand{\listtablename}{My List of Tables}


% ********************** TOC depth and numbering depth *************************

\setcounter{secnumdepth}{3}
\setcounter{tocdepth}{2}


% ******************************* Nomenclature *********************************

% To change the name of the Nomenclature section, uncomment the following line

%\renewcommand{\nomname}{Symbols}

% List-of-Abbreviations link infrastructure (bidirectional hyperlinks).
% \abbrevdef{KEY}      -- place at the FIRST use of an acronym in a chapter.
%                         (a) sets a hyperlink target named abbrev:KEY so the
%                             List of Abbreviations can link back to here, and
%                         (b) renders the acronym itself as a hyperlink to the
%                             matching row in the List of Abbreviations
%                             (target nom:KEY). Clicking jumps to the table.
% \abbrevref{KEY}{TXT} -- used in Nomenclature/nomenclature.tex to render the
%                         definition cell as a hyperlink to the first-use site.
% \nomrow{KEY}{TXT}    -- used in Nomenclature/nomenclature.tex to render one
%                         row of the abbreviation table. The acronym cell
%                         carries a hypertarget nom:KEY so the chapter
%                         occurrences can link to it.
% Made \abbrev robust so it survives "moving arguments" (TOC/LOF/LOT entries,
% \caption short forms, section bookmarks). Without \protect the embedded
% \hyperlink macro expands prematurely and corrupts the .toc/.lof/.lot files.
\DeclareRobustCommand{\abbrevdef}[1]{\hypertarget{abbrev:#1}{\hyperlink{nom:#1}{#1}}}
\DeclareRobustCommand{\abbrev}[1]{\hyperlink{nom:#1}{#1}}
\newcommand{\abbrevref}[2]{\hyperlink{abbrev:#1}{#2}}
% The anchor is placed in a zero-size box so that hyperref's raised link does not
% carry the visible acronym below the baseline of its own row: with the anchor and the
% text in one \hypertarget group, the left column rendered half a line under the right.
\newcommand{\nomrow}[2]{%
  \raisebox{0pt}[0pt][0pt]{\hypertarget{nom:#1}{}}#1 & \abbrevref{#1}{#2} \\}


% ********************************* Appendix ***********************************

% The default value of both \appendixtocname and \appendixpagename is `Appendices'. These names can all be changed via:

%\renewcommand{\appendixtocname}{List of appendices}
%\renewcommand{\appendixname}{Appndx}

% *********************** Configure Draft Mode **********************************

% Uncomment to disable figures in `draft'
%\setkeys{Gin}{draft=true}  % set draft to false to enable figures in `draft'

% These options are active only during the draft mode
% Default text is "Draft"
%\SetDraftText{DRAFT}

% Default Watermark location is top. Location (top/bottom)
%\SetDraftWMPosition{bottom}

% Draft Version - default is v1.0
%\SetDraftVersion{v1.1}

% Draft Text grayscale value (should be between 0-black and 1-white)
% Default value is 0.75
%\SetDraftGrayScale{0.8}


% ******************************** Todo Notes **********************************
%% Uncomment the following lines to have todonotes.

%\ifsetDraft
%	\usepackage[colorinlistoftodos]{todonotes}
%	\newcommand{\mynote}[1]{\todo[author=kks32,size=\small,inline,color=green!40]{#1}}
%\else
%	\newcommand{\mynote}[1]{}
%	\newcommand{\listoftodos}{}
%\fi

% Example todo: \mynote{Hey! I have a note}